
\documentclass{article}
\usepackage{colortbl}
\usepackage{makecell}
\usepackage{multirow}
\usepackage{supertabular}

\begin{document}

\newcounter{utterance}

\centering \large Interaction Transcript for game `ifeval', experiment `ifeval', episode 367 with Qwen3.5{-}9B.
\vspace{24pt}

{ \footnotesize  \setcounter{utterance}{1}
\setlength{\tabcolsep}{0pt}
\begin{supertabular}{c@{$\;$}|p{.15\linewidth}@{}p{.15\linewidth}p{.15\linewidth}p{.15\linewidth}p{.15\linewidth}p{.15\linewidth}}
    \# & \multicolumn{2}{c}{Player} && \multicolumn{2}{c}{Game Master} \\
    \hline

    \theutterance \stepcounter{utterance}  
    & & & \multicolumn{4}{p{0.6\linewidth}}{
        \cellcolor[rgb]{0.9,0.9,0.9}{
            \makecell[{{p{\linewidth}}}]{
                \texttt{\tiny{[P1$\langle$GM]}}
                \texttt{Expand the riddle into a story with a funny tone:} \\
\\ 
\texttt{What can you catch but not throw?} \\
\texttt{A cold} \\
\\ 
\texttt{Use * to highlight at least 2 sections in your text. For example: *this is a highlighted text section*.} \\
            }
        }
    }
     \\ \\

    \theutterance \stepcounter{utterance}  
    & \multicolumn{4}{p{0.6\linewidth}}{
        \cellcolor[rgb]{0.561,0.737,0.561}{
            \makecell[{{p{\linewidth}}}]{
                \texttt{\tiny{[P1$\rangle$GM]}}
                \texttt{Barnaby was a man who took the phrase "catch a cold" quite literally. While everyone else treated the common cold as a metaphorical sneeze that happens to you, Barnaby believed it was a literal projectile waiting to be intercepted. He spent his days standing in the park with a net, scanning the horizon for the elusive, invisible bug.} \\
\\ 
\texttt{One particularly chilly Tuesday, Barnaby spotted it. A tiny, shivering speck drifted past a passing pigeon.} \\
\\ 
\texttt{"Gotcha!" Barnaby shouted, lunging forward with his net. He successfully scooped up the cold, holding it triumphantly in his gloved hand. It felt surprisingly light, like holding a bag of frozen peas.} \\
\\ 
\texttt{Now came the tricky part. The riddle said you could catch it, but *you cannot throw a cold*. Barnaby frowned, looking at the wriggling packet of misery in his net. "This is a problem," he muttered. "If I can't throw it, how do I get rid of it? Do I have to keep it forever? Is this a lifetime commitment?"} \\
\\ 
\texttt{He tried tossing it into the trash can, but the cold simply bounced off the lid and landed back in his hand. He tried blowing it away, but it just multiplied, turning into three smaller colds that danced around his ears.} \\
\\ 
\texttt{Panic set in. Barnaby realized the true horror of the situation. He was now the sole custodian of a contagious disaster, unable to return it to the sender. He began to shiver uncontrollably, his nose running like a faucet, and his voice cracking into a deep, nasal baritone.} \\
\\ 
\texttt{"Oh no," Barnaby whimpered, rubbing his burning eyes. "I think I accidentally *caught a cold* in the most literal sense possible, and now I'm stuck with it."} \\
\\ 
\texttt{As he sneezed a stream of glittery germs onto his own sweater, Barnaby finally understood the riddle's punchline. You can catch a cold, yes, but the moment you do, you're stuck with it until it runs its course. There is no throwing it back, no returning it to the sky. You just have to suffer through the sniffles, the fever, and the general misery, hoping that eventually, the cold will just... go away on its own.} \\
\\ 
\texttt{"Great," Barnaby groaned, wrapping himself in a blanket. "Just great. Now I have to explain to my boss why I'm working from home with a net full of germs."} \\
            }
        }
    }
    & & \\ \\

    \theutterance \stepcounter{utterance}  
    & & & \multicolumn{2}{p{0.3\linewidth}}{
        \cellcolor[rgb]{0.9,0.9,0.9}{
            \makecell[{{p{\linewidth}}}]{
                \texttt{\tiny{[GM$|$GM]}}
                {'detectable_format:number_highlighted_sections': {'num_highlights': 2}}
            }
        }
    }
    & & \\ \\

    \theutterance \stepcounter{utterance}  
    & & & \multicolumn{2}{p{0.3\linewidth}}{
        \cellcolor[rgb]{0.9,0.9,0.9}{
            \makecell[{{p{\linewidth}}}]{
                \texttt{\tiny{[GM$|$GM]}}
                \texttt{game\_result = WIN} \\
            }
        }
    }
    & & \\ \\

\end{supertabular}
}

\end{document}
